\subsection{The integral closure of a ring. Integrally closed domains}

DEFINITION 3. Let A be a ring and R a (commutative) A-algebra. The sub-A-algebra A' of R consisting of the elements of R integral over A (no. 1, Corollary 1 to Proposition 4) is called the integral closure of A in R. If A' is equal to the canonical image of A in R, A is called integrally closed in R.

\textbf{Remarks}\par

\begin{enumerate}
\def\labelenumi{(\arabic{enumi})}
\item
  If \(h: \mathbf{A} \to \mathbf{R}\) is the ring homomorphism defining the A-algebra structure on \(\mathbf{R}\) , the integral closure of \(\mathbf{A}\) in \(\mathbf{R}\) is also that of \(h(\mathbf{A})\) in \(\mathbf{R}\) . On the other hand, if \(\mathbf{R}'\) is a subalgebra of \(\mathbf{R}\) , the integral closure of \(\mathbf{A}\) in \(\mathbf{R}'\) is \(\mathbf{A}' \cap \mathbf{R}'\) .
\item
  If A is a field, the integral closure \(A'\) of A in R consists of the elements of R which are algebraic over A (no. 1, Example 1); generalizing the terminology used for field extensions (Algebra, Chapter V, §3, no. 3), \(A'\) is then also called the algebraic closure of the field A in the algebra R and A is called algebraically closed in R if \(A' = A\) .
\end{enumerate}

DEFINITION 4. If A is an integral domain, the integral closure of A in its field of fractions is called the integral closure of A. An integral domain is called integrally closed if it is equal to its integral closure.

\pandocbounded{\includegraphics[keepaspectratio]{images/part002_6bc214f1.jpg}}

Note that an integrally closed domain is not necessarily closed in any ring containing it, as the example of a field which is not algebraically closed shows.

PROPOSITION 7. Let A be a ring and R an A-algebra. The integral closure A' of A in R is a subring integrally closed in R.

The integral closure of A' in R is integral over A by no. 1, Proposition 6; it is therefore equal to A'.

COROLLARY. The integral closure of an integral domain A is an integrally closed domain.

Let K be the field of fractions of A and B the integral closure of A. Clearly

K is the field of fractions of B and it is sufficient to apply Proposition 7 to R = K.

PROPOSITION 8. Let R be a ring, \((\mathbf{B}_{\lambda})_{\lambda \in \mathbf{L}}\) a family of subrings of R and for each \(\lambda \in L\) let \(A_{\lambda}\) be a subring of \(B_{\lambda}\) . If each \(A_{\lambda}\) is integrally closed in \(B_{\lambda}\) , then \(A = \bigcap_{\lambda \in L} A_{\lambda}\) is integrally closed in \(B = \bigcap_{\lambda \in L} B_{\lambda}\) .

This follows immediately from Definition 3 and no. 1, Corollary 1 to Proposition 2.

COROLLARY. Every intersection of a non-empty family of integrally closed subdomains of an integral domain A is an integrally closed domain.

It is sufficient to apply Proposition 8 taking R and the \(B_{\lambda}\) equal to the field of fractions K of A and noting that a subring of K integrally closed in K is a fortiori an integrally closed domain since its field of fractions is contained in K.

PROPOSITION 9. Let A be a ring, \((\mathbf{R}_i)_{1 \leqslant i \leqslant n}\) a finite family of A-algebras and \(\mathbf{A}_i'\) the integral closure of A in \(\mathbf{R}_i\) ( \(1 \leqslant i \leqslant n\) ). Then the integral closure of A in \(\mathbf{R} = \prod_{i=1}^{n} \mathbf{R}_i\) is equal to \(\prod_{i=1}^{n} \mathbf{A}_i'\) .

This is an immediate consequence of no. 1, Proposition 3.

COROLLARY 1. Let A be a reduced Noetherian ring, \(\mathfrak{p}_i\) ( \(1 \leqslant i \leqslant n\) ) its distinct minimal prime ideals, \(\mathbf{K}_i\) the field of fractions of the integral domain \(\mathbf{A} / \mathfrak{p}_i\) (canonically isomorphic to the local ring \(\mathbf{A}_{\mathfrak{p}_i}\) (Chapter IV, §2, no. 5, Proposition 10)) and \(\mathbf{A}_i'\) the integral closure of A in \(\mathbf{K}_i\) ( \(1 \leqslant i \leqslant n\) ). Then the canonical isomorphism of the total ring of fractions B of A onto \(\prod_{i=1}^{n} \mathbf{K}_i\) (loc. cit.) maps the integral closure of A in B onto the product ring \(\prod_{i=1}^{n} \mathbf{A}_i'\) .

COROLLARY 2. For a reduced Noetherian ring to be integrally closed in its ring of fractions it is necessary and sufficient that it be a direct composition of integrally closed (Noetherian) domains.

